\documentclass[11pt]{article}
\usepackage[margin=1in]{geometry}
\usepackage{booktabs}
\usepackage{tabularx}
\usepackage{multirow}
\usepackage{makecell}
\usepackage{array}

\title{Tabularx Multirow Stress Case}
\author{Test Corpus}
\date{}

\begin{document}
\maketitle

This test combines expandable columns with row spans and explicit line breaks
inside cells. The table is useful for identifying parsers that assume fixed-width
tabular environments or single-line cell content.

\begin{table}[htbp]
\centering
\caption{Program features by implementation phase}
\label{tab:tabularx}
\begin{tabularx}{\linewidth}{ll>{\raggedright\arraybackslash}X>{\centering\arraybackslash}p{0.16\linewidth}}
\toprule
Phase & Component & Implementation detail & Status \\
\midrule
\multirow{3}{*}{Design}
  & Eligibility & \makecell[l]{Sites define local thresholds\\using the shared protocol.} & Complete \\
  & Intake & Intake forms include a short baseline screen and a referral checklist. & Complete \\
  & Training & Supervisors receive a facilitation guide with worked examples. & Active \\
\midrule
\multirow{2}{*}{Delivery}
  & Monitoring & Teams submit weekly notes describing deviations from the planned sequence. & Active \\
  & Escalation & \makecell[l]{Cases crossing risk thresholds\\move to senior review.} & Pilot \\
\midrule
\multirow{2}{*}{Review}
  & Audit & Reviewers reconcile the service log with appointment records. & Pending \\
  & Feedback & Sites receive annotated summaries before the next reporting cycle. & Pending \\
\bottomrule
\end{tabularx}
\end{table}

\end{document}
